\scanpage{38}
\mspage{(3)}
\textit{Application aux corps de nombres algébriques.} --- Soit $\mathrm A$ un corps de
nombres algébriques; à toute valeur absolue $v$ de $\mathrm A$, on associe le
complété $\mathrm A_v$ de $\mathrm A$ pour cette valeur absolue. D'où une application
\[
\mathrm Q(\mathrm A)\to\prod\mathrm Q(\mathrm A_v)
\]
et cette application est injective (théorème de Hasse).

Si $v$ est une valuation discrète, toute classe de forme sur $\mathrm A_v$
est déterminée par $\varepsilon$, $w^1$ et $w^2$ $^{(4)}$.
Si $\mathrm A_v=\mathbf R$, le corps des réels, alors $\mathrm H(\mathbf R)$
s'identifie à $\mathbf Z/2\mathbf Z[t]$. Toute classe de forme est définie par
deux nombres entiers $m$ et $n$ ($m$, nombre de carrés positifs; $n$, nombre de
carrés négatifs).

Posant $n=\sum n_j2^j$ ($n_j=0$ ou $1$), on trouve
\[
w^j(q)=\binom nj\bmod 2;\qquad\text{en particulier}\quad w^{2^j}=n_j.
\]
Donc une classe de forme sur $\mathbf R$ est complètement déterminée par $w$.
De tout ce qui précède on déduit le diagramme commutatif suivant :
\diagram{\begin{tikzcd}
\mathrm Q(\mathrm A)\arrow[r]\arrow[d,"w"']&
\prod\mathrm Q(\mathrm A_v)\arrow[d,"(w_v)"]\\
\mathrm H(\mathrm A)\arrow[r]&\prod\mathrm H(\mathrm A_v)
\end{tikzcd}}
On démontre ainsi que l'application $w$ est injective : Sur un corps de
nombres toute classe de forme quadratique est déterminée par son rang et ses
classes de Stiefel-Whitney.

\textit{Remarque.} --- On ne sait pas si l'application
\[
\mathrm{grad}^2(\mathrm Q(\mathrm A))\to\mathrm H^2(\mathrm A,\mathbf Z/2\mathbf Z)
\]
est toujours injective, et l'on ne sait pas non plus si tout élément d'ordre 2
du groupe de Brauer de $\mathrm A$ est somme d'éléments de la forme $a.b$.

$^{(1)}$ N. \textsc{Bourbaki}, \textit{Algèbre}, chap. IX, Hermann, Paris, 1959.

$^{(2)}$ A. \textsc{Grothendieck}, \textit{Bull. Soc. math. Fr.}, \textbf{86}, 1958.

$^{(3)}$ T. A. \textsc{Springer}, \textit{Akad. Amst. Proc.}, \textbf{62 A}, 1959.

$^{(4)}$ E. \textsc{Witt}, \textit{Journal de Crelle}, \textbf{176}, 1937.
\mspage{37}
